\part{Proof of Concept and Delivered MVP}

\chapter{Proof of Concept}

The assignment's proof-of-concept requirement is intentionally modest: demonstrate working frontend to
backend communication, validate the core API data contract, and organize the repository according to the
architecture. The delivered system satisfies that minimum while going beyond it.

\section{PoC Minimum Expected by the Assignment}

The brief requires three core things:

\begin{enumerate}
  \item a working frontend with a visible editing page,
  \item at least one meaningful API call between frontend and backend,
  \item request/response data contracts that match the architecture document.
\end{enumerate}

The delivered repository satisfies those requirements directly:

\begin{itemize}
  \item the frontend provides a working document editor interface,
  \item the backend exposes concrete routes for document creation, loading, saving, permissions, AI,
  sessions, versioning, revert, and export,
  \item the documented contracts are represented by OpenAPI, contract notes, and automated tests.
\end{itemize}

\section{Delivered MVP Features Beyond the Minimum PoC}

The final implementation goes well beyond the minimum PoC threshold while still preserving the same
architectural skeleton. The delivered MVP includes:

\begin{itemize}
  \item a visible local login experience backed by \texttt{POST /auth/login},
  \item local bearer-token login through \texttt{POST /auth/login},
  \item live collaboration with Yjs over WebSocket,
  \item autosave with idempotent backend persistence,
  \item version history and revert,
  \item document-scoped permissions and AI policy controls,
  \item async AI jobs for rewrite, summarize, and translate,
  \item partial AI acceptance,
  \item synchronous and asynchronous export flows,
  \item Swagger/OpenAPI docs at \texttt{/docs},
  \item root commands for install, run, test, and deterministic demo reset,
  \item Docker Compose and CI support.
\end{itemize}

\section{Repository and README Evidence}

The assignment rubric explicitly checks whether the repository is organized according to the documented
structure and whether a third party can run the system from the README. The delivered repository provides
that through:

\begin{itemize}
  \item a monorepo matching the documented service decomposition,
  \item root orchestration scripts: \texttt{install:all}, \texttt{dev:all}, \texttt{test:all},
  and \texttt{dev:clean},
  \item service-specific READMEs and a root README,
  \item contract notes and OpenAPI for backend routes,
  \item a deterministic smoke test that validates backend + realtime wiring.
\end{itemize}

\section{Verification Strategy}

The delivered MVP demonstrates feasibility through both automated and manual validation.

\begin{itemize}
  \item Backend tests validate routing, auth, permissions, AI jobs, export, versions, and sessions.
  \item Frontend tests validate API clients, realtime behavior, AI panel interactions, and document-page
  workflows such as autosave and stale-save handling.
  \item Realtime tests validate websocket session behavior, collaborative convergence, viewer rejection,
  and revert propagation.
  \item A root integration smoke test validates that backend and realtime services work together under a
  shared SQLite state.
\end{itemize}

In addition, the live demo flow shows the user-facing proof that the architecture holds together:
login, create a document, grant access to a second user, live collaboration, autosave, AI invocation,
partial AI acceptance, version history, revert, export, and Swagger docs.

\section{Running the AI Flow with LM Studio}

The delivered MVP supports two AI-provider modes:

\begin{itemize}
  \item \texttt{stub}: deterministic demo-safe behavior with no external dependency,
  \item \texttt{lmstudio}: a local OpenAI-compatible provider for a real model-backed experience.
\end{itemize}

For the strongest demo, LM Studio should be started before the application stack so the backend can use
the real provider path.

Recommended LM Studio setup for the delivered MVP:

\begin{enumerate}
  \item Open LM Studio.
  \item Load a local model that supports the OpenAI-compatible chat-completions API.
  \item Start the local server and keep it reachable at \texttt{http://127.0.0.1:1234}.
  \item Verify that the model server responds at \texttt{/v1/models}.
  \item Start the application with either:
  \begin{itemize}
    \item \texttt{npm run dev:clean} (auto-detects LM Studio when available), or
    \item \texttt{AI\_PROVIDER=lmstudio npm run dev:clean} to force the LM Studio path.
  \end{itemize}
\end{enumerate}

This design is intentional. The stub provider guarantees reliable demos and deterministic tests, while
LM Studio demonstrates the same asynchronous AI workflow against a real local model endpoint without
requiring a paid hosted provider.

Practical verification commands are:

\begin{lstlisting}[language={}]
curl http://127.0.0.1:1234/v1/models
AI_PROVIDER=lmstudio npm run dev:clean
\end{lstlisting}

If LM Studio is unavailable, the rest of the system --- editing, collaboration, versioning, permissions,
and export --- still remains usable. Only the AI-generation path degrades, which is exactly the graceful
failure behavior described in the architecture.
